\documentclass[../../main.tex]{subfiles}% 注意这里的文件路径不能用 ./main.tex ,否则用latexmk编译子文件会报错
\graphicspath{{\subfix{./image/}}{\subfix{../../image/}}} % 指定图片引用路径,后续可以直接使用图片文件名
% 如果图片存放在上级目录,则使用 ../../image/ 作为备用路径

% 例如:
% \begin{figure}[H]
% \centering
% \includegraphics[width=0.3\textwidth]{图.png}
% \caption{}
% \label{figure:图}
% \end{figure}
% 注意:上述\label{}一定要放在\caption{}之后,否则引用图片序号会只会显示??.

\begin{document}

\section{存在性结果的证明：逼近问题 (Pn)}\label{section:ch06-逼近问题}

\subsection{初值的逼近}\label{subsection:ch06-初值逼近}

设 (A0) 成立。为处理问题 (P)，需要初值 $u_0$ 的合适逼近。为此，使用如下技术引理。

\begin{lemma}\label{lemma:ch06-6.5.1}
对任意 $\mu\in\mathscr{R}^+_f(\Omega)$，存在序列 $\{(\alpha_n,\beta_n,\gamma_n)\}\subseteq[C^\infty_c(\Omega)]^3$，$\alpha_n\geqslant0$，$\beta_n\geqslant0$，$\gamma_n\geqslant0$ 在 $\Omega$ 上，具有如下性质：
\begin{align}
\alpha_n\to\mu_r \text{ 在 }L^1(\Omega), \label{align:ch6-6.41a}
\end{align}
\begin{align}
\alpha_n\stackrel{*}{\rightharpoonup}\mu_{\mathrm{ac}},\quad
\beta_n\stackrel{*}{\rightharpoonup}[\mu_s]_{d,2},\quad
\gamma_n\stackrel{*}{\rightharpoonup}\mu_{c,2}
\quad\text{在 }\mathscr{R}_f(\Omega), \label{align:ch6-6.41b}
\end{align}
\begin{align}
\int_\Omega\beta_n f_n\rho\,\mathrm{d}x\to\langle[\mu_s]_{d,2},f\rho\rangle \label{align:ch6-6.41c}
\end{align}
对所有 $\rho\in C_c(\Omega)$ 及所有 $\{f_n\}\subseteq H^1(\Omega)\cap L^\infty(\Omega)$，$\sup_{n\in\mathbb{N}}\|f_n\|_\infty<\infty$，$f_n\rightharpoonup f$ 在 $H^1(\Omega)$ 成立。

置 $\mu_n\triangleq \alpha_n+\beta_n+\gamma_n\in C^\infty_c(\Omega)$。可取 $\alpha_n,\beta_n,\gamma_n$ 使
\begin{align}
\|\mu_n\|_{L^1(\Omega)}\leqslant\|\mu\|_{\mathscr{R}_f(\Omega)},\quad
\|\mu_n\|_{L^\infty(\Omega)}\leqslant\frac{M\|\mu\|_{\mathscr{R}_f(\Omega)}}{\sqrt[4]{n}}
\quad\text{对某个 }M>0, \label{align:ch6-6.41d}
\end{align}
\begin{align}
\mu_n\to\mu_r \text{ 在 }\Omega\text{ 中几乎处处},\quad
\mu_n\stackrel{*}{\rightharpoonup}\mu \text{ 在 }\mathscr{R}_f(\Omega). \label{align:ch6-6.41e}
\end{align}
\end{lemma}

\begin{proof}
置 $\widetilde\mu\triangleq \widetilde\mu_r+\widetilde\mu_s\in\mathscr{R}^+_f(\mathbb{R}^N)$，
\begin{align*}
\widetilde\mu_r(x)\triangleq \mu_r(x)\chi_\Omega(x)\quad(x\in\mathbb{R}^N),\qquad
\widetilde\mu_s(E)\triangleq \mu_s(\Omega\cap E)\quad(E\in\mathscr{B}^d).
\end{align*}
由定义 $\widetilde\mu=\widetilde\mu\llcorner\Omega$，故对每个 $E\in\mathscr{B}^d\cap\Omega$，$\widetilde\mu(E)=\mu(E)$。考虑序列 $\{\zeta_n\}\subseteq C^\infty_c(\mathbb{R}^N)$，定义为
\begin{align}
\zeta_n(x)\triangleq \frac{n^{d\theta}}{\int_{\mathbb{R}^N}\zeta(x)\,\mathrm{d}x}\zeta(n^\theta x) \quad(x\in\mathbb{R}^N), \label{align:ch6-6.42}
\end{align}
其中 $\zeta\in C^\infty_c(\mathbb{R}^N)$ 为标准磨光子，$\theta>0$ 待选。对任意 $n\in\mathbb{N}$，置
\begin{align*}
\widetilde\alpha_n\triangleq \widetilde\mu_{\mathrm{ac}}*\zeta_n,\quad
\widetilde\beta_n\triangleq [\widetilde\mu_s]_{d,2}*\zeta_n,\quad
\widetilde\gamma_n\triangleq \widetilde\mu_{c,2}*\zeta_n,\quad
\widetilde\mu_n\triangleq \widetilde\alpha_n+\widetilde\beta_n+\widetilde\gamma_n=\widetilde\mu*\zeta_n
\end{align*}
（见\eqref{align:ch03-3.82}，$p=2$，$m=1$）。设 $\Omega_n$ 为开集，$\Omega_n\Subset\Omega_{n+1}\Subset\Omega$ 对每个 $n\in\mathbb{N}$，$\bigcup_{n=1}^\infty\Omega_n=\Omega$，设 $\eta_n\in C^\infty_c(\Omega_{n+1})$，$0\leqslant\eta_n\leqslant1$，$\eta_n=1$ 在 $\Omega_n$ 上（$n\in\mathbb{N}$）。

我们证明函数
\begin{align*}
\alpha_n\triangleq \widetilde\alpha_n\eta_n,\quad
\beta_n\triangleq \widetilde\beta_n\eta_n,\quad
\gamma_n\triangleq \widetilde\gamma_n\eta_n,\quad
\mu_n\triangleq \widetilde\mu_n\eta_n=\alpha_n+\beta_n+\gamma_n\quad(n\in\mathbb{N})
\end{align*}
（它们显然在 $\Omega$ 上非负）具有\eqref{align:ch6-6.41a}中所述性质。

\eqref{align:ch6-6.41a}--\eqref{align:ch6-6.41b}中的收敛断言由标准卷积结果得到。为证\eqref{align:ch6-6.41c}，固定 $\rho\in C^1_c(\Omega)$。置 $F_n\triangleq (f_n\rho)*\zeta_n$，取 $n$ 充分大使 $\operatorname{supp}\rho\subseteq\Omega_n$ 且 $\operatorname{supp}F_n\subseteq\Omega$。注意 $\eta_n=1$ 在 $\Omega_n$ 上，得
\begin{align}
\langle[\mu_s]_{d,2},F_n\rangle
=\langle[\widetilde\mu_s]_{d,2}*\zeta_n,\eta_n f_n\rho\rangle
=\int_\Omega\beta_n f_n\rho\,\mathrm{d}x. \label{align:ch6-6.43}
\end{align}
注意由 Banach--Alaoglu 定理 $f_n\stackrel{*}{\rightharpoonup}f$ 在 $L^\infty(\Omega)$，因为由假设 $\{f_n\}$ 在 $L^\infty(\Omega)$ 中有界。显然 $F_n\stackrel{*}{\rightharpoonup}f\rho$ 在 $L^\infty(\Omega)$。事实上，对每个 $\zeta\in L^1(\Omega)$，有
\begin{align}
\begin{aligned}
\lim_{n\to\infty}\int_\Omega F_n\zeta\,\mathrm{d}x
&=\lim_{n\to\infty}\int_\Omega(f_n\rho)*\zeta_n\,\zeta\,\mathrm{d}x\\
&=\lim_{n\to\infty}\int_\Omega\zeta(x)\left(\int_\Omega f_n(y)\rho(y)\zeta_n(x-y)\,\mathrm{d}y\right)\mathrm{d}x\\
&=\lim_{n\to\infty}\int_\Omega f_n(y)\rho(y)\,\mathrm{d}y\int_\Omega\zeta(x)\zeta_n(x-y)\,\mathrm{d}x
=\int_\Omega f\rho\zeta\,\mathrm{d}x.
\end{aligned} \label{align:ch6-6.44}
\end{align}
由\eqref{align:ch6-6.44}也有 $F_n\rightharpoonup f\rho$ 在 $L^2(\Omega)$。因 $f_n\rho\in H^1_0(\Omega)$，对每个充分大的 $n$，$F_n\in H^1_0(\Omega)$，弱梯度 $\nabla F_n=[\nabla(f_n\rho)]*\zeta_n$。于是显然 $F_n\rightharpoonup f\rho$ 在 $H^1_0(\Omega)$，因为由假设 $f_n\rightharpoonup f$ 在 $H^1(\Omega)$。

因 $[\mu_s]_{d,2}\in L^1(\Omega)+H^{-1}(\Omega)$（见\cref{theorem:ch03-扩散部分H负1}），由上述收敛结果得
\begin{align*}
\lim_{n\to\infty}\langle[\mu_s]_{d,2},F_n\rangle=\langle[\mu_s]_{d,2},f\rho\rangle.
\end{align*}
在\eqref{align:ch6-6.43}中令 $n\to\infty$，即得\eqref{align:ch6-6.41c}对所有 $\rho\in C^1_c(\Omega)$，由标准逼近结果断言随之成立。

\eqref{align:ch6-6.41d}中第一个不等式由 $\mu_n$ 的定义直接得到，第二个通过在\eqref{align:ch6-6.42}中取 $\theta\in(0,1/(4N)]$ 得到：
\begin{align*}
\|\mu_n\|_{L^\infty(\Omega)}\leqslant Cn^{N\theta}\|\mu\|_{\mathscr{R}_f(\Omega)}
\leqslant\frac{C}{\sqrt[4]{n}}\|\mu\|_{\mathscr{R}_f(\Omega)}.
\end{align*}

现在处理\eqref{align:ch6-6.41e}。置
\begin{align}
[\mu_s]_n\triangleq \beta_n+\gamma_n\equiv(\widetilde\mu_s*\zeta_n)\eta_n, \label{align:ch6-6.45}
\end{align}
注意由 Radon--Nikodým 定理，对几乎所有 $x\in\Omega$，
\begin{align}
\lim_{n\to\infty}\alpha_n(x)=\mu_r(x), \label{align:ch6-6.46}
\end{align}
\begin{align}
\limsup_{n\to\infty}[\mu_s]_n(x)
\leqslant\frac{1}{\int_{\mathbb{R}^N}\zeta(x)\,\mathrm{d}x}
\lim_{n\to\infty}n^{N\theta}\widetilde\mu_s\left(B\left(x,\frac{1}{n^\theta}\right)\right)=0 \label{align:ch6-6.47}
\end{align}
（见\eqref{align:ch02-2.77}中第二个不等式，\cref{corollary:ch02-Lebesgue微分}）。因
\begin{align*}
\mu_n=\alpha_n+[\mu_s]_n,
\end{align*}
结合\eqref{align:ch6-6.46}与\eqref{align:ch6-6.47}得
\begin{align*}
\mu_n(x)\to\mu_r(x) \quad\text{对几乎所有 }x\in\Omega,
\end{align*}
\eqref{align:ch6-6.41e}中第一个收敛得证（\eqref{align:ch6-6.41e}中第二个收敛另见\eqref{align:ch6-6.41b}）。

对任意 $u_0\in\mathscr{R}_f(\Omega)$，考虑 Jordan 分解 $u_0=u^+_0-u^-_0$，其中 $u^\pm_0\in\mathscr{R}^+_f(\Omega)$，$u^+_0\perp u^-_0$。设 $u^\pm_{0r}\in L^1(\Omega)$ 为 $[u^\pm_0]_{\mathrm{ac}}$ 的密度，$\delta^\pm\triangleq [u^\pm_0]_{d,2}$，$\gamma^\pm\triangleq [u^\pm_0]_{c,2}$ 分别为 $u^\pm_0$ 关于 $2$-容量的弥散部分与集中部分。由\eqref{align:ch03-3.82}有 $u^\pm_0=\gamma^\pm+\delta^\pm$。置
\begin{align}
[u^\pm_0]_n\triangleq \gamma^\pm_n+\delta^\pm_n\geqslant0,\quad
u_{0n}\triangleq [u^+_0]_n-[u^-_0]_n, \label{align:ch6-6.48}
\end{align}
其中 $\{\gamma^\pm_n\}\subseteq C^\infty_c(\Omega)$，$\{\delta^\pm_n\}\subseteq C^\infty_c(\Omega)$，$\gamma^\pm_n\geqslant0$，$\delta^\pm_n\geqslant0$，满足
\begin{align}
\|[u^\pm_0]_n\|_{L^1(\Omega)}\leqslant\|u^\pm_0\|_{\mathscr{M}(\Omega)},\quad
\|[u^\pm_0]_n\|_{L^\infty(\Omega)}\leqslant\frac{M\|u^\pm_0\|_{\mathscr{M}(\Omega)}}{\sqrt[4]{n}}, \label{align:ch6-6.49a}
\end{align}
\begin{align}
[u^\pm_0]_n\to u^\pm_{0r} \text{ 在 }\Omega\text{ 中几乎处处}, \label{align:ch6-6.49b}
\end{align}
\begin{align}
[u^\pm_0]_n\stackrel{*}{\rightharpoonup}u^\pm_0 \text{ 在 }\mathscr{M}(\Omega), \label{align:ch6-6.49c}
\end{align}
\begin{align}
\delta^\pm_n\stackrel{*}{\rightharpoonup}[u^\pm_0]_{d,2},\quad
\gamma^\pm_n\stackrel{*}{\rightharpoonup}[u^\pm_0]_{c,2}
\quad\text{在 }\mathscr{M}(\Omega), \label{align:ch6-6.49d}
\end{align}
\begin{align}
\int_\Omega\delta^\pm_n f_n\rho\,\mathrm{d}x\to\langle[u^\pm_0]_{d,2},f\rho\rangle \label{align:ch6-6.49e}
\end{align}
对所有 $\rho\in C_c(\Omega)$ 及 $\{f_n\}\subseteq H^1(\Omega)\cap L^\infty(\Omega)$，$\sup_{n\in\mathbb{N}}\|f_n\|_{L^\infty(\Omega)}<\infty$，$f_n\rightharpoonup f$ 在 $H^1(\Omega)$ 成立。此后总假设 (Pn) 的初值 $u_{0n}$ 具有性质\eqref{align:ch6-6.49a}。
\end{proof}

\begin{remark}\label{remark:ch06-6.5.1}
对任意 $u_0\in\mathscr{R}_f(\Omega)$，考虑分解（见\eqref{align:ch03-3.82}，$p=2$，$m=1$）
\begin{align}
u^\pm_0=[u^\pm_0]_{\mathrm{ac}}+[u^\pm_{0s}]_{d,2}+[u^\pm_0]_{c,2}. \label{align:ch6-6.50}
\end{align}
对 $\mu=u^\pm_0$ 应用\cref{lemma:ch06-6.5.1}，得序列 $\{(\alpha^\pm_n,\beta^\pm_n,\gamma^\pm_n)\}\subseteq[C^\infty_c(\Omega)]^3$ 满足\eqref{align:ch6-6.41a}。置 $\delta^\pm_n\triangleq \alpha^\pm_n+\beta^\pm_n$，性质\eqref{align:ch6-6.49a}立得。
\end{remark}

\begin{lemma}\label{lemma:ch06-6.5.2}
设 $\{u_{0n}\}$ 满足\eqref{align:ch6-6.48}与\eqref{align:ch6-6.49a}--\eqref{align:ch6-6.49c}。设 $f\in C(\mathbb{R})$ 满足\cref{proposition:ch05-函数弱星}。则存在 $\{u_{0n}\}$ 的子列（不重新标号），使
\begin{align}
f\circ u_{0n}\stackrel{*}{\rightharpoonup}f\circ u_{0r}+M^+_f u^+_{0s}-M^-_f u^-_{0s}
\quad\text{在 }\mathscr{R}_f(\Omega). \label{align:ch6-6.51}
\end{align}
特别地，$u_{0n}\stackrel{*}{\rightharpoonup}u_0$ 在 $\mathscr{R}_f(\Omega)$。
\end{lemma}

\begin{proof}
设 $\{\nu_n\}\subseteq\mathcal{Y}^+(\Omega;\mathbb{R})$ 为与 $\{u_{0n}\}$ 相伴的 Young 测度序列。由\eqref{align:ch6-6.49a}中第一个不等式，序列 $\{u_{0n}\}$ 在 $L^1(\Omega)$ 中有界，故由\cref{theorem:ch05-弱紧Young}存在 Young 测度 $\nu\in\mathcal{Y}^+(\Omega;\mathbb{R})$，使（可能抽取子列，不重新标号）$\nu_n\to\nu$ 窄收敛。由\eqref{align:ch6-6.49b}中的收敛、\cref{proposition:ch05-几乎处处收敛}及等式\eqref{align:ch05-5.53}，$\nu$ 的离散分解 $\{\nu_x\}_{x\in\Omega}$ 满足对几乎所有 $x\in\Omega$，$\nu_x=\delta_{u_{0r}(x)}$。于是
\begin{align*}
\int_\mathbb{R}y^\pm\,\mathrm{d}\nu_x(y)=u^\pm_{0r}(x)
\quad\text{对几乎所有 }x\in\Omega.
\end{align*}
因此，如\cref{remark:ch05-biting注}中那样论证，存在 $\{u_{0n}\}$ 的子列（不重新标号）及 $\sigma^\pm\in\mathscr{R}^+_f(\Omega)$，使
\begin{align}
[u_{0n}]^\pm\stackrel{*}{\rightharpoonup}u^\pm_{0r}+\sigma^\pm \quad\text{在 }\mathscr{R}_f(\Omega). \label{align:ch6-6.52}
\end{align}
由 $u_{0n}$ 的定义（见\eqref{align:ch6-6.48}）及\eqref{align:ch6-6.49c}有 $u_{0n}\stackrel{*}{\rightharpoonup}u^+_0-u^-_0=u_0$，而由\eqref{align:ch6-6.52}，$u_{0n}\stackrel{*}{\rightharpoonup}u_{0r}+\sigma$，其中 $\sigma\triangleq \sigma^+-\sigma^-$。

注意由\eqref{align:ch6-6.48}中第二个等式，对所有 $n\in\mathbb{N}$ 有 $[u_{0n}]^\pm\leqslant[u^\pm_0]_n$，故由\eqref{align:ch6-6.49b}--\eqref{align:ch6-6.49c}与\eqref{align:ch6-6.52}得 $\sigma^\pm\leqslant u^\pm_{0s}$ 在 $\mathscr{R}_f(\Omega)$。这蕴含 $\sigma^\pm$ 关于 Lebesgue 测度奇异且相互奇异。于是等式 $u_0=u_{0r}+\sigma$ 及 Lebesgue 分解的唯一性给出 $\sigma=u_{0s}$，$\sigma^\pm=u^\pm_{0s}$。然后如\cref{proposition:ch05-集中测度极限}证明中那样，结论随之成立。
\end{proof}

\subsection{先验估计}\label{subsection:ch06-先验估计}

由经典结果，对任意 $n\in\mathbb{N}$，问题 (Pn) 存在唯一经典解 $u_n\in C^\infty(Q)$（例如见文献~\cite{ref63}）；且对任意 $n\in\mathbb{N}$，有
\begin{align}
\|u_n\|_{L^\infty(0,T;L^1(\Omega))}\leqslant\|u_{0n}\|_{L^1(\Omega)},\quad
\|u_n\|_{L^\infty(Q)}\leqslant\|u_{0n}\|_{L^\infty(\Omega)}. \label{align:ch6-6.53}
\end{align}
我们需要序列 $\{u_n\}$ 与 $\{\phi(u_n)\}$ 的先验估计。这由如下命题给出。

\begin{lemma}\label{lemma:ch06-6.5.3}
设 (A0)--(A1) 成立，$\{u_n\}$ 为问题 (Pn) 的解序列。则对所有 $n\in\mathbb{N}$，有
\begin{align}
\|u_n\|_{L^\infty(0,T;L^1(\Omega))}\leqslant\|u_0\|_{\mathscr{R}_f(\Omega)}, \label{align:ch6-6.54}
\end{align}
\begin{align}
\|u_n\|_{L^\infty(Q)}\leqslant\frac{M\|u_0\|_{\mathscr{R}_f(\Omega)}}{\sqrt[4]{n}}, \label{align:ch6-6.55}
\end{align}
从而
\begin{align}
\|\phi_n(u_n)-\phi(u_n)\|_{L^\infty(Q)}
\leqslant\frac{\|u_n\|_{L^\infty(Q)}}{n}
\leqslant M\|u_0\|_{\mathscr{R}_f(\Omega)}n^{-3/4}, \label{align:ch6-6.56}
\end{align}
\begin{align}
\iint_Q|\nabla[\phi(u_n)]|^2\,\mathrm{d}x\,\mathrm{d}t
+\frac{1}{n}\iint_Q\phi'(u_n)|\nabla u_n|^2\,\mathrm{d}x\,\mathrm{d}t
\leqslant2\|\phi\|_\infty\|u_0\|_{\mathscr{R}_f(\Omega)}. \label{align:ch6-6.57}
\end{align}
\end{lemma}

\begin{proof}
不等式\eqref{align:ch6-6.49a}与\eqref{align:ch6-6.53}立即给出\eqref{align:ch6-6.54}与\eqref{align:ch6-6.55}，\eqref{align:ch6-6.56}由 $\phi_n$ 的定义直接得到（回忆 $\phi_n(z)\triangleq \phi(z)+z/n$，$z\in\mathbb{R}$）。为证\eqref{align:ch6-6.57}，将 (Pn) 第一式乘以 $\phi(u_n)$ 并在 $Q$ 上积分。利用\eqref{align:ch6-6.54}，得
\begin{align*}
\begin{aligned}
&\iint_Q|\nabla[\phi(u_n)]|^2\,\mathrm{d}x\,\mathrm{d}t
+\frac{1}{n}\iint_Q\phi'(u_n)|\nabla u_n|^2\,\mathrm{d}x\,\mathrm{d}t\\
&\quad=-\iint_Q\phi(u_n)\partial_t u_n\,\mathrm{d}x\,\mathrm{d}t\\
&\quad=\int_\Omega\mathrm{d}x\int_0^{u_{0n}(x)}\phi(y)\,\mathrm{d}y
-\int_\Omega\mathrm{d}x\int_0^{u_n(x,T)}\phi(y)\,\mathrm{d}y\\
&\quad\leqslant\|\phi\|_\infty[\|u_{0n}\|_{L^1(\Omega)}+\|u_n(\cdot,T)\|_{L^1(\Omega)}]
\leqslant2\|\phi\|_\infty\|u_0\|_{\mathscr{R}_f(\Omega)}.
\end{aligned}
\end{align*}
结论随之成立。
\end{proof}

对每个 $h\in C(\mathbb{R})$，置
\begin{align}
H(y)\triangleq \int_0^y h(z)\,\mathrm{d}z \quad(y\in\mathbb{R}). \label{align:ch6-6.58}
\end{align}

\begin{lemma}\label{lemma:ch06-6.5.4}
设 (A1) 成立，$\{u_n\}$ 为问题 (Pn) 的解序列。
\begin{enumerate}[(i)]
\item 对所有 $h\in\operatorname{Lip}(\mathbb{R})$，$\zeta\in C^1([0,T];C^1_c(\Omega))$，及 $t\in(0,T]$，有
\begin{align}
\begin{aligned}
&\int_\Omega H(u_n(x,t))\zeta(x,t)\,\mathrm{d}x
+\iint_{Q_t}h'(u_n)\nabla u_n\cdot\nabla\phi_n(u_n)\zeta\,\mathrm{d}x\,\mathrm{d}s\\
&\qquad-\iint_{Q_t}H(u_n)\partial_s\zeta\,\mathrm{d}x\,\mathrm{d}s\\
&\quad=\int_\Omega H(u_{0n}(x))\zeta(x,0)\,\mathrm{d}x
-\iint_{Q_t}h(u_n)\nabla\phi_n(u_n)\cdot\nabla\zeta\,\mathrm{d}x\,\mathrm{d}s,
\end{aligned} \label{align:ch6-6.59}
\end{align}
其中 $Q_t\triangleq \Omega\times(0,t)$。若 $h(0)=0$，则\eqref{align:ch6-6.59}对所有 $\zeta\in C^1(Q)$ 成立。
\item 对每个非减 $h\in C(\mathbb{R})$，$\zeta\in C^1([0,T];C^1_c(\Omega))$，$\zeta\geqslant0$，及 $t\in(0,T]$，
\begin{align}
\begin{aligned}
&\iint_{Q_t}H(u_n)\partial_s\zeta\,\mathrm{d}x\,\mathrm{d}s
-\iint_{Q_t}h(u_n)\nabla\phi_n(u_n)\cdot\nabla\zeta\,\mathrm{d}x\,\mathrm{d}s\\
&\quad\geqslant\int_\Omega H(u_n(x,t))\zeta(x,t)\,\mathrm{d}x
-\int_\Omega H(u_{0n}(x))\zeta(x,0)\,\mathrm{d}x.
\end{aligned} \label{align:ch6-6.60}
\end{align}
若 $h(0)=0$，则\eqref{align:ch6-6.60}对所有 $\zeta\in C^1(Q)$，$\zeta\geqslant0$ 成立。
\end{enumerate}
\end{lemma}

\begin{proof}
因 $u_n$ 光滑，等式\eqref{align:ch6-6.59}由将 (Pn) 第一式乘以 $h(u_n)\zeta$（$h\in\operatorname{Lip}(\mathbb{R})$，$h(0)=0$，$\zeta\in C^1(Q)$）并在 $Q$ 上积分直接得到。此外，若 $\zeta\in C^1([0,T];C^1_c(\Omega))$，条件 $h(0)=0$ 可省略。故 (i) 得证。

因对所有 $n\in\mathbb{N}$ 有 $\nabla u_n\cdot\nabla\phi_n(u_n)\geqslant0$，由\eqref{align:ch6-6.59}对每个非减 $h\in\operatorname{Lip}(\mathbb{R})$ 得\eqref{align:ch6-6.60}。因每个非减 $h\in C(\mathbb{R})$ 可由一列非减 Lipschitz 函数局部逼近，结论随之成立。
\end{proof}

\begin{remark}\label{remark:ch06-6.5.2}
设 (A0)--(A1) 成立，$f\in C(\mathbb{R})$ 非减，$F_\phi$ 由\eqref{align:ch6-6.9}定义。在\eqref{align:ch6-6.60}中置 $g=f\circ\phi$，得不等式
\begin{align}
\begin{aligned}
&\int_\Omega F_\phi(u_n(x,t_2))\zeta(x,t_2)\,\mathrm{d}x
-\int_\Omega F_\phi(u_n(x,t_1))\zeta(x,t_1)\,\mathrm{d}x\\
&\quad\leqslant\int_{t_1}^{t_2}\int_\Omega\{F_\phi(u_n)\partial_t\zeta
-f(\phi(u_n))\nabla[\phi_n(u_n)]\cdot\nabla\zeta\}\,\mathrm{d}x\,\mathrm{d}t
\end{aligned} \label{align:ch6-6.61}
\end{align}
对所有 $t_1,t_2\in[0,T]$，$t_1<t_2$，及 $\zeta\in C^1([0,T];C^1_c(\Omega))$，$\zeta\geqslant0$ 成立。
\end{remark}

\begin{lemma}\label{lemma:ch06-6.5.5}
设 (A0)--(A1) 成立，$\{u_n\}$ 为问题 (Pn) 的解序列。则：
\begin{enumerate}[(i)]
\item 存在 $C_1=C_1(\lambda_N(\Omega),\|u_0\|_{\mathscr{R}_f(\Omega)})>0$，使对所有 $n\in\mathbb{N}$，
\begin{align}
\iint_Q\frac{|\nabla u_n|^2}{n^2}\,\mathrm{d}x\,\mathrm{d}t\leqslant\frac{C_1}{\sqrt{n}}; \label{align:ch6-6.62}
\end{align}
\item 存在 $C_2=C_2(\lambda_N(\Omega),\|\phi\|_\infty,\|u_0\|_{\mathscr{R}_f(\Omega)})>0$，使对所有 $n\in\mathbb{N}$，
\begin{align}
\iint_Q|\nabla[\phi_n(u_n)]|^2\,\mathrm{d}x\,\mathrm{d}t\leqslant C_2. \label{align:ch6-6.63}
\end{align}
\end{enumerate}
\end{lemma}

\begin{proof}
将 $\partial_t u_n=\Delta\phi_n(u_n)$ 乘以 $u_n/n$ 并在 $Q$ 上积分，直接得
\begin{align*}
\frac{1}{2n}\int_\Omega u_n^2(x,T)\,\mathrm{d}x-\frac{1}{2n}\int_\Omega u_{0n}^2(x)\,\mathrm{d}x
\leqslant-\iint_Q\frac{|\nabla u_n|^2}{n^2}\,\mathrm{d}x\,\mathrm{d}t.
\end{align*}
由此不等式及\eqref{align:ch6-6.49a}中第二个不等式，得
\begin{align*}
\iint_Q\frac{|\nabla u_n|^2}{n^2}\,\mathrm{d}x\,\mathrm{d}t
\leqslant\frac{1}{2n}\int_\Omega u_{0n}^2\,\mathrm{d}x
\leqslant\frac{M^2\lambda_N(\Omega)}{2\sqrt{n}}\|u_0\|^2_{\mathscr{R}_f(\Omega)}
=:\frac{C_1}{\sqrt{n}}.
\end{align*}
这证明了\eqref{align:ch6-6.62}。结合\eqref{align:ch6-6.57}与\eqref{align:ch6-6.62}直接得\eqref{align:ch6-6.63}，结论成立。
\end{proof}

\begin{lemma}\label{lemma:ch06-6.5.6}
设 (A1) 成立，$\{u_n\}$ 为问题 (Pn) 的解序列。设 $h\in\operatorname{Lip}(\mathbb{R})$ 有界且非减，$h(0)=0$，$H$ 为\eqref{align:ch6-6.58}中的函数。则对任意 $\rho\in C^1_c(\Omega)$，存在 $M_{h,\rho}>0$，使对所有 $n\in\mathbb{N}$，
\begin{align}
\int_0^T\left|\frac{\mathrm{d}}{\mathrm{d}t}\int_\Omega H(u_n(x,t))\rho\,\mathrm{d}x\right|\mathrm{d}t\leqslant M_{h,\rho}. \label{align:ch6-6.64}
\end{align}
\end{lemma}

\begin{proof}
对任意 $\rho\in C^1_c(\Omega)$，置
\begin{align*}
S_{n,\rho}(t)\triangleq \int_\Omega H(u_n(x,t))\rho(x)\,\mathrm{d}x \quad(t\in(0,T)).
\end{align*}
在\eqref{align:ch6-6.59}中取 $\zeta(x,t)=\beta(t)\rho(x)$，$\beta\in C^1_c(0,T)$，表明 $S_{n,\rho}$ 有弱导数
\begin{align}
\begin{aligned}
S'_{n,\rho}(t)
={}&-\int_\Omega h(u_n(x,t))\nabla\phi_n(u_n(x,t))\cdot\nabla\rho(x)\,\mathrm{d}x\\
&-\int_\Omega h'(u_n(x,t))\nabla u_n(x,t)\cdot\nabla\phi_n(u_n(x,t))\rho(x)\,\mathrm{d}x
\end{aligned} \label{align:ch6-6.65}
\end{align}
对几乎所有 $t\in(0,T)$。由\eqref{align:ch6-6.63}，
\begin{align}
\iint_Q|h(u_n)\nabla\phi_n(u_n)\cdot\nabla\rho|\,\mathrm{d}x\,\mathrm{d}t
\leqslant\|h\|_{L^\infty(\mathbb{R})}\|\rho\|_{C^1(\overline\Omega)}
\iint_Q|\nabla\phi_n(u_n)|\,\mathrm{d}x\,\mathrm{d}t
\leqslant C_2\|h\|_{L^\infty(\mathbb{R})}\|\rho\|_{C^1(\overline\Omega)}. \label{align:ch6-6.66}
\end{align}
估计\eqref{align:ch6-6.65}右端第二项。在\eqref{align:ch6-6.59}中取 $\zeta=1$，$t=T$，利用\eqref{align:ch6-6.49a}中第一个不等式，直接得
\begin{align}
\begin{aligned}
0\leqslant\iint_Q h'(u_n)\nabla u_n\cdot\nabla\phi_n(u_n)\,\mathrm{d}x\,\mathrm{d}t
\leqslant\int_\Omega H(u_{0n})\,\mathrm{d}x
\leqslant\|h\|_{L^\infty(\mathbb{R})}\|u_{0n}\|_{L^1(\Omega)}
\leqslant\|h\|_{L^\infty(\mathbb{R})}\|u_0\|_{\mathscr{R}_f(\Omega)}
\end{aligned} \label{align:ch6-6.67}
\end{align}
（注意 $H(u_n(\cdot,T))\geqslant0$ 几乎处处在 $\Omega$ 上，因 $h$ 非减且 $h(0)=0$）。由\eqref{align:ch6-6.65}--\eqref{align:ch6-6.67}有 $S'_{n,\rho}\in L^1(0,T)$，且
\begin{align}
\int_0^T|S'_{n,\rho}(t)|\,\mathrm{d}t
\leqslant(C_2\|\rho\|_{C^1(\overline\Omega)}+\|\rho\|_{C(\overline\Omega)}\|u_0\|_{\mathscr{R}_f(\Omega)})
\|h\|_{L^\infty(\mathbb{R})}T. \label{align:ch6-6.68}
\end{align}
结论随之成立。
\end{proof}

\begin{lemma}\label{lemma:ch06-6.5.7}
设 (A0)--(A1) 成立，$\{u_n\}$ 为问题 (Pn) 的解序列。则：
\begin{enumerate}[(i)]
\item 对每个 $t\in(0,T)$，
\begin{align}
\int_\Omega|\nabla[\phi_n(u_n)]|^2(x,s)\,\mathrm{d}x\leqslant\frac{C_2}{t}
\quad\text{对所有 }s\in(t,T]\text{ 及 }n\in\mathbb{N} \label{align:ch6-6.69}
\end{align}
其中 $C_2>0$ 由\cref{lemma:ch06-6.5.5}(ii) 给出；
\item 对所有 $t\in(0,T)$ 及 $n\in\mathbb{N}$，
\begin{align}
\int_t^T\int_\Omega[\phi'_n(u_n)](\partial_\tau u_n)^2\,\mathrm{d}x\,\mathrm{d}\tau
\leqslant\frac{C_2}{2t}. \label{align:ch6-6.70a}
\end{align}
进而，若 (A2) 成立，则对每个 $t\in(0,T)$，
\begin{align}
\int_t^T\int_\Omega(\partial_\tau[\phi_n(u_n)])^2\,\mathrm{d}x\,\mathrm{d}\tau
\leqslant\frac{C_2(1+\|\phi'\|_\infty)}{2t}. \label{align:ch6-6.70b}
\end{align}
\end{enumerate}
\end{lemma}

\begin{proof}
(i) 对任意固定 $t\in(0,T)$，置 $h(s)\triangleq \frac{s}{t}\chi_{[0,t)}(s)+\chi_{[t,T)}(s)$（$s\in(t,T]$）。将 $\partial_t u_n=\Delta\phi_n(u_n)$ 乘以 $h(t)\partial_t[\phi_n(u_n)]$ 并在 $Q_s=\Omega\times(0,s)$ 上积分，对每个 $s\in(t,T)$，得
\begin{align}
\begin{aligned}
0&\leqslant\iint_{Q_s}h(\tau)\phi'_n(u_n)(\partial_\tau u_n)^2\,\mathrm{d}x\,\mathrm{d}\tau\\
&=-\int_0^s\mathrm{d}\tau\,h(\tau)\int_\Omega\partial_\tau\{\nabla[\phi_n(u_n)]\}\cdot\nabla[\phi_n(u_n)]\,\mathrm{d}x\\
&=-\frac{1}{2}\int_0^s h(\tau)\frac{\mathrm{d}}{\mathrm{d}\tau}\left(\int_\Omega|\nabla[\phi_n(u_n)]|^2\,\mathrm{d}x\right)\mathrm{d}\tau\\
&=-\frac{1}{2}\int_\Omega|\nabla[\phi_n(u_n)]|^2(x,s)\,\mathrm{d}x
+\frac{1}{2t}\int_0^t\mathrm{d}\tau\int_\Omega|\nabla[\phi_n(u_n)]|^2\,\mathrm{d}x.
\end{aligned} \label{align:ch6-6.71}
\end{align}
利用\eqref{align:ch6-6.63}，由上述不等式得\eqref{align:ch6-6.69}。故 (i) 得证。

(ii) 在\eqref{align:ch6-6.71}中取 $s=T$，得
\begin{align*}
\int_t^T\int_\Omega[\phi'_n(u_n)](\partial_\tau u_n)^2\,\mathrm{d}x\,\mathrm{d}\tau
\leqslant\iint_Q h(\tau)\phi'_n(u_n)(\partial_\tau u_n)^2\,\mathrm{d}x\,\mathrm{d}\tau
\leqslant\frac{1}{2t}\int_0^t\mathrm{d}\tau\int_\Omega|\nabla[\phi_n(u_n)]|^2\,\mathrm{d}x
\leqslant\frac{C_2}{2t}.
\end{align*}
这证明了\eqref{align:ch6-6.70a}，\eqref{align:ch6-6.70b}直接成立。证明完成。
\end{proof}

\begin{remark}\label{remark:ch06-6.5.3}
因
\begin{align*}
|\nabla[\phi_n(u_n)]|^2=\left(\phi'(u_n)+\frac{1}{n}\right)^2|\nabla u_n|^2
\geqslant(\phi'(u_n))^2|\nabla u_n|^2=|\nabla\phi(u_n)|^2,
\end{align*}
由\eqref{align:ch6-6.69}，对所有 $0<t<s\leqslant T$，得
\begin{align}
\int_\Omega|\nabla\phi(u_n)|^2(x,s)\,\mathrm{d}x\leqslant\frac{C_2}{t}. \label{align:ch6-6.72}
\end{align}
类似地，因 $|\partial_\tau\phi(u_n)|^2\leqslant|\partial_\tau[\phi_n(u_n)]|^2$，由\eqref{align:ch6-6.70b}得
\begin{align}
\int_t^T\int_\Omega(\partial_\tau[\phi(u_n)])^2\,\mathrm{d}x\,\mathrm{d}\tau
\leqslant\frac{C_2(1+\|\phi'\|_\infty)}{2t}. \label{align:ch6-6.73}
\end{align}
\end{remark}

\begin{lemma}\label{lemma:ch06-6.5.8}
设 (A0)--(A1) 成立，$\{u_n\}$ 为问题 (Pn) 的解序列。则：
\begin{enumerate}[(i)]
\item 对所有 $k>0$ 及 $n\in\mathbb{N}$，有
\begin{align}
\iint_{\{\pm u_n>k\}}|\nabla\phi(u_n)|^2\,\mathrm{d}x\,\mathrm{d}t
\leqslant\|u^\pm_0\|_{\mathscr{R}_f(\Omega)}[\pm\phi_{\pm\infty}\mp\phi(\pm k)], \label{align:ch6-6.74}
\end{align}
\item 存在 $M_0>0$，使对所有 $k>0$ 及 $n\in\mathbb{N}$，
\begin{align}
\iint_{\{\pm u_n>k\}}|\nabla\phi_n(u_n)|^2\,\mathrm{d}x\,\mathrm{d}t
\leqslant M_0[\pm\phi_{\pm\infty}\mp\phi(\pm k)+n^{-1/2}]. \label{align:ch6-6.75}
\end{align}
\end{enumerate}
\end{lemma}

\begin{proof}
我们仅证\eqref{align:ch6-6.74}中上号情形，另一情形类似。对任意 $k>0$，置
\begin{align*}
F(y)\triangleq \chi_{(k,+\infty)}(y)\int_k^y[\phi(z)-\phi(k)]\,\mathrm{d}z \quad(y\in\mathbb{R}).
\end{align*}
于是 $F'\in W^{1,\infty}_{\mathrm{loc}}(\mathbb{R})$，且
\begin{align}
0\leqslant F(y)\leqslant[\phi_{+\infty}-\phi(k)](y-k)\chi_{(k,+\infty)}(y). \label{align:ch6-6.76}
\end{align}
将 $\partial_t u_n=\Delta\phi_n(u_n)$ 乘以 $F'(u_n)$ 并在 $Q$ 上积分，由\eqref{align:ch6-6.76}与\eqref{align:ch6-6.49a}得
\begin{align*}
\begin{aligned}
&\iint_{\{u_n>k\}}|\nabla\phi(u_n)|^2\,\mathrm{d}x\,\mathrm{d}t
+\frac{1}{n}\iint_{\{u_n>k\}}\phi'(u_n)|\nabla u_n|^2\,\mathrm{d}x\,\mathrm{d}t\\
&\quad=\int_\Omega F(u_{0n}(x))\,\mathrm{d}x-\int_\Omega F(u_n(x,T))\,\mathrm{d}x\\
&\quad\leqslant[\phi_{+\infty}-\phi(k)]\int_{\{u_{0n}\geqslant k\}}[u_{0n}(x)-k]\,\mathrm{d}x
\leqslant\|u^+_0\|_{\mathscr{R}_f(\Omega)}[\phi_{+\infty}-\phi(k)].
\end{aligned}
\end{align*}
于是\eqref{align:ch6-6.74}得证。关于\eqref{align:ch6-6.75}，由\eqref{align:ch6-6.74}与\eqref{align:ch6-6.62}有
\begin{align*}
\begin{aligned}
\iint_{\{\pm u_n>k\}}|\nabla\phi_n(u_n)|^2\,\mathrm{d}x\,\mathrm{d}t
&\leqslant2\iint_{\{\pm u_n>k\}}\left(|\nabla\phi(u_n)|^2+\frac{|\nabla u_n|^2}{n^2}\right)\mathrm{d}x\,\mathrm{d}t\\
&\leqslant2\|u_0\|_{\mathscr{R}_f(\Omega)}[\pm\phi_{\pm\infty}\mp\phi(\pm k)]
+2C_1n^{-1/2}
\end{aligned}
\end{align*}
其中 $C_1>0$ 如\cref{lemma:ch06-6.5.5}(i)。置 $M_0\triangleq 2\max\{\|u_0\|_{\mathscr{R}_f(\Omega)},C_1\}$，得\eqref{align:ch6-6.75}，结论成立。
\end{proof}




\end{document}